\documentclass[11pt,a4paper]{article}

\usepackage[margin=2.35cm]{geometry}
\usepackage{amsmath}
\usepackage{booktabs}
\usepackage{enumitem}
\usepackage{siunitx}
\usepackage{hyperref}

\hypersetup{colorlinks=true,linkcolor=blue,urlcolor=blue,citecolor=blue}
\setlist{nosep}

\title{Trinkey Sampling, BNO055 Timing, and Compensator Evaluation}
\author{}
\date{}

\begin{document}

\maketitle

\section{Recommended decision}

Use a nominal paired-sensor acquisition rate of \SI{250}{Hz} with the BNO055
accelerometer bandwidth set to \SI{125}{Hz} for the planned \SI{10}{Hz} and
\SI{35}{Hz} print comparisons.  Retain Klipper's \SI{1}{kHz} controller grid.
These rates form a valid multirate system: the command may be generated at
\SI{1}{kHz} while measured acceleration arrives every \SI{4}{ms}.

For the present evaluation, do not invent the three intervening acceleration
values between sensor samples.  Compare the desired acceleration with the
measured acceleration at the genuine, time-mapped BNO055 sample times.  This
gives a valid \emph{acceleration-domain} tracking metric.  It does not reproduce
the paper's position-domain RMS tracking error unless an observer or another
position estimate is added.

The two principal evaluation outputs should therefore be:

\begin{enumerate}
    \item printed-part evidence: photographs and repeated dimensional
    measurements; and
    \item a time-aligned acceleration-domain comparison from Motan, reported as
    an adaptation necessitated by the available instrumentation.
\end{enumerate}

This is a defensible dissertation limitation.  It avoids claiming that
interpolated data contain information that the sensors did not measure.

\section{What the current acquisition actually measures}

The streaming record contains separate timestamps for the base/bed and
toolhead accelerometers.  The firmware reads the base sensor and then the
toolhead sensor over I\textsuperscript{2}C.  For each sensor, the RP2040
timestamp is recorded \emph{after} the register read succeeds.  The host then
maps that timestamp from the Trinkey clock to Klipper print time using the
fitted affine clock relation

\begin{equation}
    t_{\mathrm{print}} = \alpha t_{\mathrm{Trinkey}} + \beta .
\end{equation}

This design is better than assigning one shared timestamp to both sensors, and
it removes host USB arrival time from the signal alignment.  However, the
timestamp is not the instant at which the BNO055 internally sampled the MEMS
element or updated its output register.  The BNO055 does not provide a
per-sample timestamp register.  Consequently, the remaining timing uncertainty
includes the age of the register value, the sequential read time, and the
sensor's internal filter delay.  Exact data-ready timing would require each
BNO055 interrupt/data-ready output to be wired to a separate RP2040 GPIO and
timestamped at the edge.  Even then, the edge would identify data availability,
not the exact physical centre of the filtered measurement~\cite{bno055}.

Given the remaining project time, retaining the present hardware and reporting
this uncertainty is preferable to changing the accelerometers.  The same
sensor and timing path is used in every controller condition, so repeatable
latency is largely common to the comparisons, although it must still be treated
consistently when computing phase-sensitive errors.

\section{Why the requested \SI{400}{Hz} stream was not \SI{400}{Hz}}

The reported status at the end of the supplied run contained
\(45{,}438\) scheduled sample indices and \(18{,}266\) missed indices.  Thus
only

\begin{equation}
    45{,}438-18{,}266=27{,}172
\end{equation}

paired acquisitions were performed, or approximately \SI{60}{\percent} of the
requested samples.  The approximately \SI{3.8}{ms} worst acquisition loop
cannot meet the \SI{2.5}{ms} period required by \SI{400}{Hz}.  The absence of
transport, host-queue, packet, and frame errors shows that the bottleneck was
the sequential sensor acquisition rather than USB transfer or Motan.

At \SI{250}{Hz}, the available period becomes \SI{4}{ms}, which is only about
\SI{0.2}{ms} above the observed worst loop time.  Therefore \SI{250}{Hz} is a
candidate operating point, not an assumption.  Before accepting a run, verify
that the following counters do not increase throughout a complete print:

\begin{equation}
    \mathtt{firmware\_dropped}=
    \mathtt{firmware\_overruns}=
    \mathtt{firmware\_transport\_dropped}=0,
\end{equation}

and also require zero host queue drops, packet-sequence errors, frame errors,
and read errors.  A batch size of four concerns transport only; it does not
mean that four samples share one timestamp.

The \SI{250}{Hz} stream provides 25 samples per cycle at \SI{10}{Hz} and about
7.1 samples per cycle at \SI{35}{Hz}.  This is adequate for the intended
fundamental-frequency comparison.  It would provide only 2.5 samples per cycle
at \SI{100}{Hz}; Nyquist is satisfied, but the margin is weak and broadband
results near \SI{100}{Hz} should not be presented as high-fidelity
measurements.  The \SI{125}{Hz} BNO055 bandwidth setting also leaves limited
anti-alias margin at a \SI{250}{Hz} acquisition rate.  For the present tests,
frequency-selective results at \SI{10}{Hz} and \SI{35}{Hz} are more defensible
than a full-band RMS claim extending to \SI{100}{Hz}.

\section{Evaluation procedure}

The generated trajectories provide matched uncompensated and filtered
B-spline prints at \SI{10}{Hz} and \SI{35}{Hz}.  Their motion and extrusion
payloads are identical at a given frequency; only the controller state changes.
This supports a controlled two-by-two comparison.

\subsection{Dynamic metric from Motan}

The Motan toolhead TrapQ is the nominal desired trajectory.  The filtered
B-spline wrapper modifies the stepper kinematic position without replacing
that nominal TrapQ.  Consequently, the reference should be
\(\mathtt{trapq(toolhead,x\_accel)}\) or
\(\mathtt{trapq(toolhead,y\_accel)}\), and the measured output should be:

\begin{itemize}
    \item toolhead accelerometer X for X-axis motion; and
    \item bed/base accelerometer Y for Y-axis motion.
\end{itemize}

For each genuine accelerometer timestamp \(t_k\), evaluate the corresponding
TrapQ acceleration directly at \(t_k\).  The lightweight capture endpoint does
this inside Klipper and stores the result beside the physical sample as
\(\mathtt{trinkey\_accel(toolhead,command\_x)}\) or
\(\mathtt{trinkey\_accel(base,command\_y)}\).  It therefore avoids exporting
the complete high-density TrapQ.  This is preferable to first creating two
independently rounded \SI{250}{Hz} sequences.  If a uniform grid is needed by
an analysis tool, resample both signals once onto the same \SI{250}{Hz} grid
after clock mapping and document the interpolation and filtering used.

After applying the measured sensor orientation/sign, scale and static-bias
correction, define

\begin{equation}
    e_a[k] = a_{\mathrm{ref}}(t_k)-a_{\mathrm{meas}}(t_k),
    \qquad
    \operatorname{RMSE}_a =
    \sqrt{\frac{1}{N}\sum_{k=1}^{N} e_a[k]^2}.
    \label{eq:accel-rmse}
\end{equation}

Report Equation~\eqref{eq:accel-rmse} in \si{mm.s^{-2}}, optionally normalised
by the RMS reference acceleration.  Also fit the \SI{10}{Hz} or \SI{35}{Hz}
sine and cosine components to report error amplitude and phase.  This
narrow-band result is less sensitive to out-of-band noise and aliasing than an
unqualified broadband RMS value.

Use the same steady central-wall interval in every run and exclude warm-up,
parking and the real end-cap corner transients.  State explicitly whether a
constant instrumentation delay was corrected.  If no delay correction is
made, call the result \emph{end-to-end acceleration error}; this remains useful
for relative compensated/uncompensated comparison because the acquisition path
is common, but delay contributes to the numerical error.

Do not label Equation~\eqref{eq:accel-rmse} as the paper's position tracking
error.  For a nearly single-frequency error only, an equivalent displacement
phasor can be estimated from

\begin{equation}
    X_e(j\omega)=-\frac{A_e(j\omega)}{\omega^2},
    \qquad \omega=2\pi f .
\end{equation}

This conversion must be described as narrow-band; double-integrating a general
accelerometer record will otherwise amplify bias and drift.  A genuine
position-domain RMS comparison requires the Luenberger observer used in the
reference study, an equivalently validated estimator, or an independent
position measurement.

\subsection{Printed-part metric}

The proposed physical evaluation follows the reference paper's two forms of
evidence: visible print quality and a repeated geometric measurement.  For each
frequency and controller condition:

\begin{enumerate}
    \item keep filament, temperature, layer height, flow, acceleration limits,
    thermal state and all other settings fixed; keep input shaping disabled;
    \item photograph the same central wall region with fixed lighting,
    orientation, distance and scale;
    \item measure the nominal \SI{3.4}{mm} outside width at ten predefined,
    equally spaced locations in the central region; and
    \item report the mean width, standard deviation, mean signed dimensional
    error, and the measurement locations.  Repeat whole prints where time
    permits so print-to-print variability is not confused with controller
    performance.
\end{enumerate}

The paper compared uncompensated, standard FBF and hybrid FBF cases at low and
high excitation frequencies, using both part measurements and RMS tracking
error.  The present uncompensated-versus-filtered-B-spline comparison therefore
follows the same evaluation structure, but with fewer controller conditions and
an acceleration-domain dynamic metric until a position observer is
available~\cite{chou2022}.

\section{Implication for the later hybrid controller}

A \SI{1}{kHz}, \(140\)-sample controller window spans \SI{140}{ms}.  A
\SI{250}{Hz} sensor supplies only 35 independent measurement epochs in that
window.  This is not a problem for the evaluation above: its error sum simply
uses the 35 real epochs.  It \emph{is} an implementation issue for the hybrid
controller if its lifted model expects 140 output values.

The appropriate future solution is a multirate observer: predict the estimated
state every \SI{1}{ms} from the known command and apply a measurement correction
whenever a new \SI{250}{Hz} accelerometer sample arrives.  The observer then
supplies a \SI{1}{kHz} estimated output sequence to the lifted controller while
preserving the fact that only every fourth step contains new measurement
information.  Simple interpolation or sample repetition may produce a vector
of length 140, but does not create 140 independent observations and should not
be presented as such.  Timestamped measurement delay must also be included in
the observer or batch-history alignment.

Thus the final decision is:

\begin{itemize}
    \item \textbf{baseline compensator evaluation now:} use the real
    \SI{250}{Hz} samples, without missing-sample estimation;
    \item \textbf{paper-comparable position error:} add a validated observer;
    and
    \item \textbf{future hybrid feedback implementation:} retain the
    \SI{1}{kHz} controller and use an explicit multirate observer/residual
    construction rather than treating interpolated samples as measurements.
\end{itemize}

\begin{thebibliography}{9}

\bibitem{chou2022}
C.-H. Chou, M. Duan, and C. E. Okwudire,
\newblock ``A Physics-Guided Data-Driven Feedforward Tracking Controller for
Systems with Unmodeled Dynamics---Applied to 3D Printing,'' 2022.
\newblock doi: \href{https://doi.org/10.48550/arXiv.2206.11960}%
{10.48550/arXiv.2206.11960}.

\bibitem{bno055}
Bosch Sensortec,
\newblock ``BNO055 Intelligent 9-axis absolute orientation sensor: Data sheet,''
BST-BNO055-DS000.
\newblock \url{https://www.bosch-sensortec.com/media/boschsensortec/downloads/datasheets/bst-bno055-ds000.pdf}.

\end{thebibliography}

\end{document}
